%% kapa: exam=1 round=31 year=2020 period=2 form=A subject=appraisal_law subject_ko=감정평가관계법규 seq=2 number=2
%% source: Q-Net 감정평가사 기출문제 (cst00309) artlSeq=5207219 file=extracted/1008/5207219/A형/A형/2020년도 제31회 감정평가사 1차 시험_2교시_A형.pdf
%% conversion: kordoc 4.24.51 -> tools/gen1.py (auto); review=proofread-sonnet; flags=none
\begin{question}{2}
국토의 계획 및 이용에 관한 법령상 도시ㆍ군기본계획에 관한 설명으로 옳지 않은 것은?
\begin{choices}
  \item 다른 법률에 따른 지역ㆍ지구 등의 지정으로 인하여 도시ㆍ군기본계획의 변경이 필요한 경우에는 토지적성평가를 하지 아니할 수 있다.
  \item 광역시장은 도시ㆍ군기본계획을 변경하려면 관계 행정기관의 장과 협의한 후 지방도시계획위원회의 심의를 거쳐야 한다.
  \item 시장 또는 군수는 도시ㆍ군기본계획을 변경하려면 도지사의 승인을 받아야 한다.
  \item 시장 또는 군수는 10년마다 관할 구역의 도시ㆍ군기본계획에 대하여 그 타당성 여부를 전반적으로 재검토하여 정비하여야 한다.
  \item 「수도권정비계획법」에 의한 수도권에 속하지 아니하고 광역시와 경계를 같이하지 아니한 시로서 인구 10만명 이하인 시의 시장은 도시기본계획을 수립하지 아니할 수 있다.
\end{choices}
\end{question}
\answer{4}
